\section{PDEs}
We require PDEs for the heavy species densities, the electric
potential, and the electron distribution function.  In this initial
work, the background gas temperature $T_g$ is taken as given and thus
no gas energy equation is necessary.

\subsection{Species Continuity}
The species transport equations are as follows:
%
\begin{equation}
\pp{n_{\alpha}}{t} + \pp{F_{\alpha}}{x} = \dot{\omega}_{\alpha},
\label{eqn:species}
\end{equation}
%
where $t$ is time, $x$ is the coordinate in the electrode-normal
direction, $n_{\alpha}$ denotes the number density of species
$\alpha$, $F_{\alpha}$ is the corresponding flux, and
$\dot{\omega}_{\alpha}$ is the net production rate of species $\alpha$
due to chemical reactions.  The flux is closed as follows:
%
\begin{equation*}
F_{\alpha} = \mu_{\alpha} n_{\alpha} E - D_{\alpha} \pp{n_{\alpha}}{x},
\end{equation*}
%
where $z_{\alpha}$ is the charge number, $\mu_{\alpha}$ is the
mobility, $D_{\alpha}$ is the diffusion coefficient, and $E$ is the
electric field.  The transport closures are as given in
\S\ref{sec:fluid-closures}.  The electric field is determined from the
potential, $E = -\pp{\phi}{x}$, which is governed by a Poisson
equation (\S\ref{sec:hybrid-poisson}).  Finally, the heavy species
production due to chemical reactions are discussed in
\S\ref{sec:hybrid-chemistry}.

For a mixture of $N_s$ species, $N_s - 2$ species densities are
evolved according to~\eqref{eqn:species}.  The electron density is
obtained from the electron distribution function, which is governed by
the electron Boltzmann equation (\S\ref{sec:hybrid-ebolt}).
Further,the density of the dominant background species (e.g., argon)
is determined from the pressure, which is assumed known, and the ideal
gas law:
%
\begin{equation*}
p = \sum_{\alpha = 1}^{N_s} n_{\alpha} k_B T_{\alpha},
\end{equation*}
%
where $k_B$ is the Boltzmann constant and $T_{\alpha}$ is the
temperature of species $\alpha$.  A two temperature model is adopted
where the heavy species share a single temperature $T_g$, which is
different from the electron temperature, such that the ideal gas law
becomes
%
\begin{equation*}
p = n_e k_B T_e + k_B T_g \sum_{\alpha \neq e} n_{\alpha}.
\end{equation*}
%
The gas temperature is assumed known.  The electron temperture is
governed by the electron Boltzmann equation (\S\ref{sec:ebolt}).

\subsection{Electromagnetics} \label{sec:hybrid-poisson}
A quasi-static model is used for the electromagnetics.  In this case,
Maxwell's equations reduce to a Poisson equation governing the
electric potential $\phi$:
%
\begin{equation*}
-\pp{^2 \phi}{x^2} = \frac{\rho_c}{\epsilon_0}
\end{equation*}
%
where $\rho_c$ is the charge density and $\epsilon_0$ is the
permittivity of free space.  The charge density is given from the
species densities:
%
\begin{equation*}
\rho_c = \sum_{\alpha=1}^{N_s} z_{\alpha} q_e n_{\alpha},
\end{equation*}
%
where $q_e$ is the elementary charge (the magnitude of the charge of
an electron).

The electron number density $n_e$ is determined from the electron
distribution function,
%
\begin{equation*}
n_e(x,t) = \int_{\mbb{R}^3} f_e(x, \mbf{v}_e, t) \, d^3 \mbf{v}_e
\end{equation*}
%
which is governed by the electron Boltzmann equation.

\subsection{Electron Boltzmann Equation} \label{sec:hybrid-ebolt}
Let $f_e$ denote the electron distribution function.  In the 1-D glow
discharge model, the distribution function depends on space $x \in
(0,h)$, velocity $\mbf{v} \in \mathbb{R}^3$, and time $t \in
\mathbb{R}^+$.  Since $f_e$ depends only on $x$ and $\mbf{E} = E
\hat{x}$, the electron Boltzmann equation simplifies to
%
\begin{equation*}
\pp{f_e}{t} + v_x \pp{f_e}{x} - \frac{q_e}{m_e} E \pp{f_e}{v_x} = S_e + C_e,
\end{equation*}
%
where $S_e+C_e$ denotes the collision term, with $S_e$ representing elastic
collisions and $C_e$ representing reactive collisions.

For $S_e$, we consider only two-body interactions.  Then, suppressing
the dependence on $x$ and $t$, $S_e$ can be written as
%
\begin{equation*}
S_e(\mbf{v}_e)
=
\sum_{\alpha=1}^{N_s}
\int \left(
f_e(\mbf{v}_e^{\prime}) f_{\alpha}(\mbf{v}_{\alpha}^{\prime}
-
f_e(\mbf{v}_e) f_{\alpha}(\mbf{v}_{\alpha})
\right) \,
W_{e \alpha}
d^3 \mbf{v}_{\alpha} d^3 \mbf{v}_e^{\prime} d^3 \mbf{v}_{\alpha}^{\prime},
\end{equation*}
%
where $W_{e \alpha}$ is the transition probability associated with the
nonreactive collision between and electron and species $\alpha$.  In
practice, we will begin by considering only non-reactive collisions
between electrons and the background gas (argon), such that $S_e$ becomes
%
\begin{equation*}
S_e (\mbf{v}_e)
=
\int \left(
f_e(\mbf{v}_e^{\prime}) f_{b}(\mbf{v}_{b}^{\prime})
-
f_e(\mbf{v}_e) f_{b}(\mbf{v}_{b})
\right) \,
W_{e b}
d^3 \mbf{v}_{b} d^3 \mbf{v}_e^{\prime} d^3 \mbf{v}_{b}^{\prime},
\end{equation*}
%
where subscript $b$ denotes the background species.  The background
species distribution function $f_b$ is assumed given by the Maxwellian
consistent with its number density $n_b$ and the gas temperature
$T_g$.  Thus,
%
\begin{equation*}
f_b(x, \mbf{v}, t)
=
n_b(x,t)
\left( \frac{m}{2 \pi k_B T_g} \right)^{3/2} 
\exp \left( \frac{-m_b \|\mbf{v}\|^2}{2 k_B T_g} \right).
\end{equation*}
%
The integral in $S_e$ may be simplified by recognizing that many
interactions considered in the integral are impossible due to
conservation of momentum and energy.  Specifically, for elastic
collisions, conservation of momentum and energy require that
%
\begin{equation*}
m_e \mbf{v}_e + m_b \mbf{v}_b
=
m_e \mbf{v}_e^{\prime} + m_b \mbf{v}_b^{\prime}
\quad
\frac{1}{2} m_e \mbf{v}_e \cdot \mbf{v}_e +
\frac{1}{2} m_b \mbf{v}_b \cdot \mbf{v}_b
=
\frac{1}{2} m_e \mbf{v}_e^{\prime} \cdot \mbf{v}_e^{\prime} + 
\frac{1}{2} m_b \mbf{v}_b^{\prime} \cdot \mbf{v}_b^{\prime}
\end{equation*}
%
Then, given the direction of the post-collision velocity
$\hat{\mbf{e}}$, the post-collision velocities may be written as
%
\begin{equation*}
PUT IN
\end{equation*}
%
and the nine-dimensional integral above may be re-written as a
five-dimensional integral:
%
\begin{equation*}
\int_{\mbb{R}^3} \int_{S^2} 
\left(
f_e(\mbf{v}_e^{\prime}) f_{b}(\mbf{v}_{b}^{\prime})
-
f_e(\mbf{v}_e) f_{b}(\mbf{v}_{b})
\right) \,
\|\mbf{v}_e - \mbf{v}_b\| \sigma_{eb}(\|\mbf{v}_e - \mbf{v}_b\|, \hat{\mbf{e}})
d\hat{\mbf{e}} \, d^3 \mbf{v}_{b},
\end{equation*}
%
where $\sigma_{eb}$ denotes the differential cross section.

To write the contribution of reactive collisions, we follow the
formulation of Ern and Giovangigli~\cite{}.  The full reaction term is
written as follows:
%
\begin{equation*}
C_e = \sum_{r=1}^{N_r} C_e^r,
\end{equation*}
%
where $N_r$ is the number of reactions and $C_e^r$ is the contributin
of the $r$th reaction.  To formulate $C_e^r$, consider the following
reaction:
%
\begin{equation*}
\sum_{j \in R^r} \chi_j
\ce{<=>}
\sum_{k \in P^r} \chi_k,
\end{equation*}
%
where $\chi_j$ denotes the $j$th species, $R^r$ is the set of
reactants, and $P^r$ is the set of products.  For reactions that
include the same species more than once, $R^r$ and $P^r$ include that
species multiple times.  For example, the set of products of the argon
ionization reaction $\ce{Ar + e <=> Ar^+ + e + e}$ is given by $P^r =
\{\ce{Ar^+}, \ce{e}, \ce{e}\}$.  Further, for a given species $i$, let
$R^r_i$ and $P^r_i$ denote the set of reactants and products for
reaction $r$ with species $i$ removed only once.  Thus, for the argon
ionization reaction example, $P^r_e = \{\ce{Ar^+}, \ce{e}\}$.

Equivalently, one may write the generic chemical reaction as
%
\begin{equation*}
\sum_{j =1}^{N_s} \nu^{r}_j\chi_j
\ce{<=>}
\sum_{k=1}^{N_s} \nu^{\prime,r}_k \chi_k,
\end{equation*}
%
where $\nu^r_{j}$ and $\nu^{\prime,r}_j$ denote the stoichiometric
coefficients of the $j$th species in the reactants and products,
respectively.  Using all of this notation, the contribution of the
$r$th reaction to the reactive collision term is given by
%
\begin{align*}
C_e^r
=&
\nu_e^r \int \left( \prod_{k \in P^r} \beta_k f_k - \prod_{j \in R^r} \beta_j f_j \right)
\frac{W_{R^r P^r}}{\prod_{j \in R^r} \beta_j} \,
\prod_{j \in R^r_e} d^3 \mbf{v}_j \, \prod_{k \in P^r} d^3 \mbf{v}_k\\
&+
\nu_e^{\prime,r} \int \left( \prod_{j \in R^r} \beta_j f_j - \prod_{k \in P^r} \beta_k f_k \right)
\frac{W_{R^r P^r}}{\prod_{j \in R^r} \beta_j} \,
\prod_{j \in R^r} d^3 \mbf{v}_j \, \prod_{k \in P^r_e} d^3 \mbf{v}_k,
\end{align*}
%
where $W_{R^r P^r}$ is the transition probability for the forward
reaction,
%
\begin{equation*}
\beta_j = \frac{h^3_{\mathrm{P}}}{a_j m_j^3},
\end{equation*}
%
and $a_j$ is the degeneracy.

For binary reactions (e.g., $\chi_0 + \chi_1 \ce{<=>} \chi_2 +
\chi_3$), this may be rewritten in terms of the differential collision
cross section.  For example, the contribution of such a reaction to
species 0 is given by
%
\begin{equation*}
C_0 = \int_{\mbb{R}^3} \int_{S^2} \left( \frac{\beta_2 \beta_3}{\beta_0 \beta_1} f_2 f_3 - f_0 f_1\right) g \sigma(g, \hat{\mbf{e}}) d\hat{\mbf{e}} d^3 \mbf{v}_1
\end{equation*}
%

Formally, reactions involving three
particles (e.g., direct ionization and recombination: $\ce{Ar + e <=>
  Ar^+ + e + e}$ are supported by this framework, but this introduces
complexity into the formulation of $W$.  While this can likely be
overcome,
\todo{Alexeev introduces a collision cross section for such reactions, but I don't understand his notation.}
 we may also avoid it by decomposing such three particle
interactions into two steps.  For instance, instead of including
direct impact ionization in the mechanism, we have
%
\begin{gather*}
\ce{Ar + e <=> Ar^{\dagger} + e} \\
\ce{Ar^{\dagger} -> Ar^+ + e},
\end{gather*}
%
where $Ar^{\dagger}$ is an intermediate species.  For these reactions,
letting $\ce{A}$ be species 1, $\ce{Ar^{\dagger}}$ be species 2, and
$\ce{Ar^+}$ be species 3, the reactive collision terms in the electron
Boltzmann equation reads
%
\begin{align*}
C_e(\mbf{v}_e)
&=
\int_{\mbb{R}^3} \int_{S^2}
\left( \frac{\beta_2}{\beta_1} f_e(\mbf{v}^{\prime}_e) f_2(\mbf{v}_2) - f_e(\mbf{v}_e) f_1(\mbf{v}_1) \right) W_{10,20} d^3 \mbf{v}_1 d^3 \mbf{v}^{\prime}_e d^3 \mbf{v}_2
\\
&+\int_{\mbb{R}^3} \int_{S^2}
\left( f_e(\mbf{v}^{\prime}_e) f_1(\mbf{v}_1) - \frac{\beta_2}{\beta_1} f_e(\mbf{v}_e) f_2(\mbf{v}_2) \right) W_{10,20} d^3 \mbf{v}^{\prime}_e d^3 \mbf{v}_1 d^3 \mbf{v}_2
+
\int_{S^2} \frac{1}{\tau} f_2(\mbf{v}_2) d\hat{\mbf{e}}.
\end{align*}
%
